%% ============================================================
%% Chapter 16: Compression After Expansion
%% ============================================================

\chapter{Compression After Expansion}
\label{ch:ch16-compression-after-expansion}
\section{The Illusion of Effortless Form}

Chapter~\ref{ch:ch15-compression-fallacy} established that description length does not identify distinguishing capacity, and that a short representation's brevity is frequently evidence of accumulated background machinery rather than proximity to an informational floor. This chapter asks where that machinery comes from, and what happens to the observer, and to the person exercising a compressed competence, once the process that built it is no longer visible.

There is a persistent pattern in the way human beings perceive completed artifacts. A finished building, a published proof, a released piece of software, a polished essay, each presents itself as a coherent, apparently simple object whose internal structure invites little reflection on the process by which it came to be. Observers encountering such objects tend, with some consistency, to draw inferences about the difficulty of producing them that track the apparent simplicity of the result rather than the generative complexity of the process. This inference is often badly calibrated. It is not, however, a universal law: some finished objects really were easy to make, and surface simplicity sometimes does correlate with low production cost. The claim developed here is about a common and consequential bias, not an exceptionless regularity.

Compression is a constitutive feature of representation: the point of producing a final form is often precisely to eliminate from view the intermediate steps, failed attempts, and structural discoveries that preceded it. A proof that displayed every exploratory conjecture, failed lemma, and notational revision through which it was produced would be unreadable. But it can be epistemically misleading when the compressed form is mistaken for the primary form, and when the prior expansion is forgotten or invisible.

\section{Expansion, Compression, and a Qualified Formal Model}

A minimal formal vocabulary states a narrow, conditional version of the asymmetry between construction and compression. Let a system $\mathcal{S}$ be a configuration of components subject to a constraint set $\mathcal{C}$; a configuration $s$ is coherent if $s \models \mathcal{C}$. Let $\Omega$ denote the space of possible configurations, and $\Omega_{\mathrm{valid}} = \{s \in \Omega : s \models \mathcal{C}\}$. Cardinality is the wrong instrument for comparing this set to $\Omega$ whenever $\Omega$ is infinite, so a measure $\mu$ appropriate to the system is fixed instead.

\begin{assumption}[Sparse validity]
For the systems considered in this chapter, $\mu(\Omega_{\mathrm{valid}})/\mu(\Omega) \ll 1$.
\end{assumption}

This is not asserted as a law governing all constrained systems; it is a condition satisfied by many of the domains discussed here, working software, load-bearing structures, valid proofs, and the results that follow are conditioned on it holding. Expansion is a process $E$ acting on an initial configuration, generating a sequence of transformations, trials, or variations; it is epistemic as much as generative, revealing structure in $\Omega_{\mathrm{valid}}$ through interaction with $\mathcal{C}$. Expansion need not be performed by the agent who later benefits from it; it can be carried out by another person, an institution, a culture, or an earlier generation of a system, and its results transmitted rather than rediscovered. Compression is a triple $(C,D,\sim)$, an encoder, decoder, and task-relevant equivalence relation such that $D(C(s)) \sim s$, possibly subject to bounded distortion; it eliminates redundancy while preserving task-relevant structure, encoding invariants discovered during expansion and suppressing intermediate derivations not required to reconstruct behavior up to $\sim$.

\begin{proposition}[Conditional dependency of compression]
\label{prop:compression-requires-expansion}
In an initially unknown task domain, constructing a reliable representation $(C,D,\sim)$ generally requires information about which variations preserve $\sim$-relevant validity. Absent an external source of this information, prior search by the agent, observation, instruction, or an inherited body of structural knowledge from elsewhere, such information must be acquired through exploration of $\Omega$ near the boundary of $\Omega_{\mathrm{valid}}$.
\end{proposition}

A compression that removes components while preserving $\sim$ requires knowing which components are invariant under $\sim$ and which are incidental; where this knowledge is not supplied externally, it can only be established by observing which variations preserve validity and which do not. The practical corollary is narrower than a universal claim: one cannot, without access to some prior exploration, one's own or someone else's, produce a reliable compressed representation of a domain one does not yet understand.

This licenses a limiting case worth naming directly, because it clarifies a distinction the rest of this chapter depends on: between an artifact's observable fluency and the developmental process, if any, that an agent using that artifact has itself undergone. Consider an agent whose entire competence in a domain consists of an inherited representation, where the agent performed none of the exploration that produced it. Large language models are the clean paradigm case. A model's fluent, low-latency completions are a compression over regularities discovered through the effortful, iterative, error-correcting production of the human text on which it was trained, genuine expansion, performed by many agents over a long history. The model's forward pass exhibits the surface signature of relegated performance, speed, fluency, the absence of visible deliberation, but nothing in that forward pass corresponds to an agent's own trajectory from deliberate monitoring to earned stabilization, because no such trajectory occurred for the system producing the output. Call this \emph{inherited compression}: a representation transmitted to an agent that performed none of the exploration behind it, as opposed to relegation, where the agent that now executes automatically is the same agent that once explored deliberately.

A related but distinct case is \emph{conditioning}: stabilization that is genuinely the agent's own, but whose triggering exposure was authored by someone else, as when repeated pairing engineers a consumer's automatic response to a brand. The consumer possesses the attentional machinery relegation requires, but the stabilization was not produced by deliberate practice at solving a task of the consumer's own. Treating fast, fluent output as evidence of relegated processing therefore risks a category error: it mistakes a property of an artifact's presentation for a property of the process that generated it. The relevant questions are whose trajectory produced the underlying stabilization, by what mechanism, and for whose objective, the agent's own task in genuine relegation, another agent's prior task transmitted rather than repeated in inherited compression, or a third party's objective imposed on the agent's own stabilization machinery in conditioning.

\section{A Probabilistic Asymmetry Between Construction and Disruption}

\begin{proposition}[Perturbation asymmetry under sparse validity]
\label{prop:asymmetry}
Suppose sparse validity holds, and let perturbations of a valid state be drawn from a distribution over $\Omega$ that is not concentrated on validity-preserving transformations. Then a randomly sampled perturbation of a valid configuration is more likely to leave $\Omega_{\mathrm{valid}}$ than to remain within it.
\end{proposition}

This is a genuine probabilistic consequence of sparse validity together with a stated assumption about the perturbation distribution; it is not, by itself, a claim that evaluation is easier than construction, and the two should not be run together. Detecting that a constraint has failed can itself be expensive, requiring measurement, proof, or diagnosis that is far from trivial. Breaking a system, detecting that it has broken, explaining why, and repairing it are four distinct operations with different cost profiles. What the proposition does support is narrower and still useful: under sparse validity and a perturbation distribution not aimed at preserving coherence, maintaining a system in $\Omega_{\mathrm{valid}}$ requires selecting from a comparatively small set of validity-preserving operations, while an unstructured perturbation is likely to exit that set.

Because compression encodes a system's behavior only up to $\sim$, compressed representations are fragile in a specific sense: perturbations that move a configuration outside the region the encoder was calibrated against can produce invalid decodings, and interpretability of the compressed form depends on the receiver's prior familiarity with the domain. This grounds a further point about evaluative ease. Agents whose primary interaction with a system is evaluative, critics, reviewers, inspectors, users, typically interact with a system already inside $\Omega_{\mathrm{valid}}$, and their task is comparatively local: identify some perturbation that would move the system into $\Omega_{\mathrm{invalid}}$. Under sparse validity, this local task samples from a different, and typically larger, region of possibility space than the task of constructing or maintaining the system in the first place. The apparent ease of finding a flaw is therefore not, by itself, evidence that avoiding that flaw during construction was equally easy.

\section{Skill as Internalized Structure: Aspect Relegation Theory}

The standard account of expertise as accumulated knowledge and trained performance is underspecified in a way that limits its ability to explain some of the most striking features of expert cognition: the disappearance of felt effort, the compression of explanation, and the gap between expert and novice perception of the same situation.

\begin{definition}[Task, known aspects, active attention]
Let a task $T = \{a_1,\ldots,a_n\}$ be a finite set of aspects, each corresponding to a constraint, sub-operation, or dependency relevant to successful execution. Let $K(t) \subseteq T$ be the known aspects the agent has represented at time $t$, whether or not currently monitored. Let the active attention set $A(t) \subseteq K(t)$ be the aspects under conscious monitoring at time $t$, subject to a bounded capacity $|A(t)| \leq k$.
\end{definition}

An aspect $a_i$ is stabilized to level $1-\varepsilon$ under policy $\pi$ and environment class $\mathcal{E}$ if
\[
\Pr(\text{constraint satisfied} \mid a_i,\pi,\mathcal{E}) \geq 1-\varepsilon:
\]
stabilization is probabilistic and context-relative rather than binary, which is why expert automaticity transfers unevenly across contexts. The attention update rule is
\[
A(t+1) = \bigl(A(t) \setminus R_t\bigr) \cup N_t \cup Q_t,
\]
where $R_t \subseteq A(t)$ is the set of aspects that have met their stabilization threshold and are relegated out of active attention, $N_t$ is a set of newly attended higher-order aspects that become tractable once capacity is freed by $R_t$, and $Q_t$ is a set of aspects, possibly previously relegated, reactivated by uncertainty, novelty, unfamiliar context, or detected error. Learning is not monotonic under this rule: as lower-level operations stabilize, new higher-order aspects can enter via $N_t$, and previously stabilized routines can be recruited back via $Q_t$ when they fail outside their training distribution. Skilled performance typically retains nonzero attention at some level, goal selection, anomaly monitoring, strategic coordination, and is better modeled as a change in what is monitored than as the disappearance of monitoring altogether.

The dual-process literature distinguishes rapid, autonomous, low-working-memory-load processing from processing that supports hypothetical reasoning and loads working memory. A substantial and well-supported class of expert Type~1 performance consists of procedures that once required controlled attention, were represented in $A(t)$ and governed by explicit monitoring, and became increasingly autonomous through practice, chunking, and probabilistic stabilization. This class does not exhaust Type~1 cognition, which also includes innate, perceptual, affective, associative, and implicitly acquired processes with no comparable history of explicit control; automatic higher-cognitive processes do not arise exclusively from a history of conscious skill acquisition, and motor-learning research on explicit and implicit training finds that both routes can lead toward automatic, dual-task-robust performance, with some paradigms actually favoring implicit acquisition. Where an expert Type~1 performance does belong to the relegation class, what appears as immediate intuition is, for that performance specifically, the residue of prior deliberation; the claim does not generalize to Type~1 cognition as a category.

A candidate neurophysiological correlate, transient hypofrontality in flow states, is worth stating as a contested hypothesis rather than an established signature. If flow involves reduced explicit supervision of well-stabilized sub-operations, some forms of flow may exhibit reduced activity in neural systems associated with self-monitoring, alongside preserved or increased activity in task-relevant control systems. Experimental work complicates a simple prefrontal-withdrawal story considerably: flow has been found associated with decreased activity in the medial prefrontal cortex and amygdala but increased activity in the inferior frontal gyri, the anterior insula, the basal ganglia, and the midbrain, a pattern of selective reallocation among frontal regions rather than uniform prefrontal withdrawal, and contemporary reviews emphasize network-level reorganization over any single mechanism of blanket deactivation. Reduced deliberative supervision can arise because supervision is no longer needed, a stabilized aspect continues to satisfy its constraints unattended, or because supervision is impaired while it is still needed, as under acute stress; these two conditions should not be treated as producing indistinguishable neural signatures, since acute-stress effects on executive function vary with severity, timing, and task in ways the flow literature does not report.

\section{Borrowed Intuition: Provenance and Formation}

The relegation model developed above describes competence built through an agent's own trajectory. A companion analysis is needed for competence that arrives already compressed, because the person exercising a judgment need not be the person whose exploration produced it. Call the resulting capacity \emph{borrowed intuition}: rules of thumb, diagnostic categories, precedents, checklists, and trained models can transmit the residue of earlier search without transmitting the search itself.

Two dimensions must be kept distinct. Provenance asks whether a compression is earned, its selection and stabilization depending on the receiving agent's own trajectory, or transmitted, the relevant selection having already occurred elsewhere and reaching the receiver in compressed form. Formation asks whether the process by which a competence is acquired is self-directed, the receiving agent controlling which situations to investigate and which objectives govern revision, or conditioned, another agent or institution structuring the receiver's exposures, feedback, or permissible actions. These dimensions are independent: an externally conditioned process can still produce earned compression if the receiver acts, observes consequences, and stabilizes a policy from those consequences, while self-directed study can remain largely transmitted if the agent chooses which authorities to consult but does not reproduce or test the inquiry underlying their conclusions.

\begin{definition}[Borrowed intuition]
A competence $C_i$ is borrowed to the extent that its present economy depends upon transmitted compression, and the receiver can exercise some practical advantage of $C_i$ without reproducing the generating trajectories.
\end{definition}

On this definition, nearly all mature human competence contains some borrowed component; the concept is not intended to divide agents into independent and dependent knowers, but to identify a variable ordinarily hidden by successful performance. Borrowed intuition can enter at several stages of a trajectory: the distinctions available, the subset activated by a situation, the policy governing movement through the activated space, or the relegated dispositions themselves. A diagnostic heuristic does not merely enlarge what an agent knows; it can make certain features salient, demote other possibilities, and change the order in which hypotheses are tested, structurally different from adding a proposition to what the agent knows.

Transformation of inherited compression through later experience requires a further, careful distinction. \emph{Synthetic compression} is the process by which an agent applies, tests, combines, revises, or re-stabilizes transmitted compression through evidence generated in its own subsequent trajectory; it is a process rather than a fifth provenance class, describing movement within the taxonomy rather than another cell beside earned and transmitted. A procedure does not become earned merely because an agent has performed it many times: repetition under invariant conditions may strengthen a disposition while providing little evidence about why it works or whether an alternative would perform better. A post-acquisition trajectory is discriminating only when it changes the receiver's relative support over the candidate policies compatible with the inherited instruction; experience that merely confirms what every candidate policy predicts may stabilize behavior without clarifying its basis. This yields a chain of non-implications:
\[
\text{temporary conditioning} \not\Rightarrow \text{retained experience} \not\Rightarrow \text{stabilized compression}.
\]

\begin{proposition}[Provenance is not erased by synthesis]
If a transmitted policy undergoes synthetic compression through a receiver's later trajectory, the resulting competence may acquire earned support without ceasing to depend historically upon transmitted structure.
\end{proposition}

Unless the later trajectory independently reconstructs every distinction and selection materially contributed by the original transmitted policy, the resulting competence remains causally dependent upon external compression. Synthetic development can alter the degree and significance of that dependence, but it does not make the originating contribution disappear. Borrowed intuition is therefore neither a permanent deficiency nor a stage every competence must leave behind; a mature competence may be deeply owned, critically understood, and extensively validated while remaining partly borrowed, which is likely the normal condition of expertise.

\section{Pretrained Systems as a Limiting Case}

A frozen pretrained language model provides an unusually clean case of transmitted compression. Its parameters are selected through optimization over a corpus generated by trajectories external to the model; the resulting system can produce fluent continuations, classifications, and explanations without having participated in the investigations, conversations, and acts of composition from which the corpus arose. A description of an experiment is not the experiment; a record of a disagreement is not participation in the disagreement. The model nevertheless benefits from those histories: regularities distributed across the corpus become available through a compressed parameterization,
\[
\text{external exploration} \to \text{recorded corpus} \to \text{parameter compression} \to \text{present competence}.
\]
A purely pretrained model with frozen parameters, no persistent memory, and no post-training interaction is a limiting case of borrowed intuition. Almost all of its stable competence is transmitted in provenance and conditioned in formation. This claim does not depend on a verdict about consciousness; it follows from the causal organization of the training pipeline. Even if one attributed understanding or some form of subjectivity to the resulting system, the provenance question would remain: a capacity can be real while the experience that made it economical belongs elsewhere.

The limiting case becomes less clean once post-training and deployment are included. Instruction tuning and reinforcement learning from human feedback alter a pretrained policy using demonstrations, preferences, and evaluations selected by additional human processes, introducing new trajectories into the model's formation without automatically making the resulting compression self-directed or earned; from the perspective of provenance, this remains largely conditioned, since humans construct prompts, rank outputs, and determine the objective through which the model is updated. Whether it also counts as earned compression depends on the boundary used to identify the agent: if the training process and the deployed model are treated as one temporally continuous system, post-training contains limited cycles of action, feedback, and policy revision; if each frozen checkpoint is treated as a distinct receiver, the deployed model inherits the result of post-training just as it inherits pretraining.

This distinction is especially important for in-context learning. A model can alter its responses when examples or instructions are supplied in the prompt, without gradient updates to its parameters. Such adaptation may be substantial within the active context, but it ordinarily disappears when that context is removed, and is therefore better described as temporary conditioning unless the context or its effects are retained. Representing a system at time $t$ as $M_t = (\theta_t, m_t, \chi_t)$, persistent parameters, persistent external memory, and current context, the system has a retained post-acquisition trajectory only if an observation can alter a later state, $(\theta_{t+1}, m_{t+1}) \neq (\theta_t, m_t)$, through a causal update attributable to that observation. A change confined to $\chi_t$ can modify present behavior without becoming persistent experience: the system adapted, but it did not retain the adaptation. This is the exact formal counterpart of the state-versus-context distinction Chapter~\ref{ch:ch10-clio-architecture} draws between operative state $\sigma$ and documentary history $H$: a documentary transition that extends $\chi_t$ without touching $(\theta_t, m_t)$ is precisely CLIO's identity transition on operative state.

\section{Skill as Compressed Structure}

Relegated aspects encode structural invariants discovered through prior practice, and the felt ease of expert performance corresponds to a small active attention set relative to the full task, $|A(t)| \ll |T|$, rather than to any reduction in the task's underlying complexity. That a task feels easy to an expert is evidence that the expert has performed enough exploration to stabilize its recurring structure, not evidence that the task is simple. Relegation introduces a partial loss of reconstructability: once aspects have left $A(t)$, intermediate steps may no longer be explicitly representable, and expert knowledge is often operationally accessible but descriptively compressed, the expert can perform the task while struggling to narrate how. This explains why expert advice is so often correct for the expert and useless for the novice: the relegated aspects are no longer visible to the advisor, so the advice skips steps that are not yet stabilized for the listener.

In commercial and media contexts this becomes more troubling, because compressed outputs are frequently presented as though no prior exploration were required to reach them, when the audience's lack of background is precisely what the compression depends on to look effortless. There is a useful distinction between compression that is reconstructible by a prepared receiver and compression that is parasitic on undisclosed prior knowledge, what might be called \emph{epistemic outsourcing}: achieving an appearance of simplicity not by encoding structure efficiently but by silently delegating the exploratory work to a receiver who is not equipped to perform it and is not told they must.

\section{Fundamental Attribution Error and the Erasure of Process}

This misperception is amplified by a well-established feature of social cognition. The fundamental attribution error describes a systematic tendency to attribute observed outcomes to dispositional properties of agents, talent, intelligence, character, while underweighting situational and processual factors. Applied to skill and production, a successful outcome tends to be attributed to the producer's intrinsic qualities rather than to the structural conditions and accumulated exploration that made it possible. This pattern is not irrational given the available evidence, outcomes are visible, processes typically are not, but the environment is incomplete, and the resulting attributions inherit that incompleteness.

A related distortion operates at the level of population inference. In many domains of skilled production, outcome distributions are heavily right-skewed: a small number of participants achieve highly visible success, and a large majority do not achieve comparable recognition regardless of comparable effort. Canonization narratives constructed retrospectively tend to read as stories of difficult genius eventually rewarded, a reconstruction that can suppress the contingency and selection effects that separated the canonized case from equally capable contemporaries whose recognition never arrived.

The asymmetry formalized in Proposition~\ref{prop:asymmetry} has a direct bearing on this pattern. Where the apparent ease of locating a flaw is read as evidence that avoiding the flaw during construction was equally easy, the inference overreaches what the proposition actually supports, since detecting or explaining a failure is not always cheap. What the formal result does support is narrower: construction and evaluation typically sample from different, and often differently sized, regions of the space of possible operations, so evaluative ease is not on its own strong evidence about constructive difficulty, in either direction.

\section{Labor and the Political Economy of Delegated Constraint Management}

The preceding analysis converges on an asymmetry in how different categories of labor tend to be perceived and rewarded. Essential work, infrastructural, repeatable, stabilizing labor such as maintenance, care, teaching, and construction, often manifests as the absence of failure rather than the presence of a discrete success, and can go relatively unnoticed even when it sustains the conditions under which other, more visible production becomes possible. Wealth confers, among other things, the capacity to delegate constraint management: supply chains, maintenance schedules, and coordination problems can be handled by others and rendered invisible to the person delegating them. This delegation is not inherently problematic, but its epistemic consequence is a tendency rather than a certainty: delegation can produce constraint blindness, a reduced capacity to perceive the structural conditions that make an apparent solution non-trivial, and where it does occur, it tends to produce a familiar genre of compressed advice, ``just do this,'' that presents a compressed form presupposing an exploration the listener must still perform but has not been told about.

\section{The Ethics of Compression}

The strong standard, that an ethically compliant compressed representation must preserve, in principle, the reconstructibility of its actual generative pathway, is too demanding and in some cases incoherent: a finished mathematical proof need not preserve the psychological history of its discovery. Three distinct standards should be kept apart. Logical reconstructibility asks whether a prepared receiver could, in principle, derive or verify the compressed content from stated premises, independent of how it was actually produced. Practical teachability asks whether the compressed form gives an unprepared receiver a workable path toward the competence it represents. Historical provenance asks whether the compressed form accurately represents the actual process, whose labor, how much exploration, over what period, that produced it.

Requiring all three uniformly is too strong; mathematics regularly and legitimately satisfies logical reconstructibility while providing neither teachability nor provenance, and this is not an ethical defect in a proof.

\begin{definition}[Non-deceptive compression]
A compressed representation is non-deceptive if it does not affirmatively misrepresent the exploration, labor, or prerequisites that produced it, whether or not it also satisfies logical reconstructibility or practical teachability.
\end{definition}

Three working criteria follow from this narrower standard: non-deception about labor, representations should not imply, where it is false, that no meaningful prior exploration was required; appropriate matching of standard to context, a representation should be evaluated against the standard it actually claims to meet, not against all three at once; and accuracy of effort distribution, representations should not systematically imply that the observed outcome is more typical of the underlying attempt distribution than it is.

\section{Conclusion}

Under a stated sparse-validity assumption and a stated assumption about perturbation distributions, maintaining coherence in a system requires selecting from a comparatively narrow set of operations, while an unstructured perturbation is likely to leave that set by an easier route; this is a conditional probabilistic result, not a universal entropy law. In human skill, a substantial and well-evidenced class of expert automaticity arises from procedures that were once deliberately controlled and became increasingly autonomous through practice, though this class does not exhaust automatic cognition. Observers who encounter only a compressed result, without its generative history, are consequently prone, not certain, but prone, to underestimate the exploration that produced it. And where the compressed result reaches an agent who never performed that exploration at all, borrowed intuition names the resulting gap between operational competence and its formative history precisely, without treating the borrowing as either counterfeit or as a deficiency every competence must eventually shed.

This is the same asymmetry Chapter~\ref{ch:ch15-compression-fallacy} located in Forth's word libraries and in trained models' latent codes, examined here from the side of the observer and the borrower rather than the side of the representation itself: a short description is evidence of accumulated machinery, and a fluent performance is evidence of a stabilization history, but in neither case does the visible artifact report whose history it was, or how much of it remains available for the receiver to inspect, revise, or trust.
